\scanpage{155}
% archive foliation: 154

Digressions combinatoires (pour Réflexions \unsure{H.4|A.4})

\uline{Chap 1} Géométrie des ens de cardinal $\leq4$

\margin{\uline{NB} Faire \unsure{polygônes} des sommets, et des côtés
\unsure{rigoureusement}...}

\uline{Chap 2} L'icosaèdre gauche : ses trois
\deleted{\unsure{géométriques}}\inserted{perspectives}

\deleted{Orientations \unsure{de} $\unsure{\omega(S)}$}
$\alpha(\unsure{F})\in\omega(S)$ défini par les trois perspectives.

\uline{Chap 3} Biicosaèdres : les trois perspectives.
\margin{\deleted{$\unsure{\mathrm{Bic}(S)}$}
\[
\mathrm{Ic}(S)\simeq
\underset{\substack{\Vert\\\mathrm{Bic}(S)}}{S^*}\times\omega(S)
\]}

\uline{Chap 4} Dualité des hexagrammes

\margin{par \uline{bihexagones}}

1. L'autodualité $\mathrm{Ens}_6\simeq\mathrm{Ens}_6$
\inserted{décrits} via
$\operatorname{Aut}(S)\xrightarrow{\sim}\operatorname{Aut}(S^*)$

Application aux groupes d'autom. des
biicosaèdres \deleted{resp.}\inserted{et} icos.

2. Structure biicosaédrale par coloriage des arêtes

3. L'autodualité (trois \uline{perspectives})
\[
\left.\begin{array}{ll}
\text{a)}&\forall s\in S,\quad s'\in S^*\qquad
 S\setminus\{s\}\simeq S^*-\{s'\}\\
\text{b)}&R\unsure{\subset|\in}\deleted{\unsure{\mathcal P_2}}
 \operatorname{Ar}(S)\times\operatorname{Ar}(S^*)\\
\text{c)}&\mathcal P_{3,3}(S)\simeq\mathcal P_{3,3}(S^*)
\end{array}\right\}\ \text{avec axiomes}
\]

4\deleted{\unsure{3}}. Le dictionnaire de la dualité
\[
\begin{array}{lll}
\text{A)}&\mathrm{Ens}_5\deleted{\unsure{\mathrm{Ens}}}
 \inserted{\simeq\mathrm{Ens}(1,5)}&\approxeq\mathrm{Biicos.}\\
&\mathrm{Ens}_5^+\simeq\mathrm{Ens}^+(1,5)&\approxeq\mathrm{Icos}\\[1ex]
\text{B)}&\mathrm{Ens}_2\times\mathrm{Ens}_4
 &\approxeq\mathrm{Ens}_{2,2,2}\approxeq\mathrm{Cub}_3\\[1ex]
\text{C)}&\underbrace{\mathrm{Ens}_{3,3}}_{\text{ditrièdres}}
 &\approxeq\mathrm{Ens}_{3,3}\quad\text{dualité}
\end{array}
\]

\margin{\deleted{$\begin{gathered}
\text{A')}\ \mathrm{Bipent}\approxeq\mathrm{Biicos.}+\text{sommet}\\
\mathrm{Pent}\approxeq\mathrm{Icos.}+\unsure{\text{sommet}}\\
\text{B')}
\end{gathered}$}}

\margin{Notion biditrièdres $\equiv$}

% The long left relocation arrow rises toward item 3.
\diagram{\begin{tikzcd}[column sep=small,row sep=large]
{} & {} & {}\\
\left[\begin{gathered}
\text{A')}\quad\mathrm{Biicos.}+\text{sommet}
\approxeq\deleted{\unsure{\mathrm{Ens}_{1,5}}\simeq\mathrm{Ens}_5}
\inserted{\unsure{\mathrm{Bipent}_5}}\\
\mathrm{Icos}+\text{---}\ \overset{\approx}{\longrightarrow}
\deleted{\mathrm{Ens}_{1,5}^{+}\simeq\mathrm{Ens}_5^{+}}
\mathrm{Pent}_5\\[2ex]
\text{B')}\quad\mathrm{Biicos}+\text{ar.}\approxeq
\deleted{\unsure{\mathrm{Ens}_2}}\times\mathrm{Car}\\[3ex]
\text{C')}\quad\mathrm{Biicos}+\text{biface}\simeq
\mathrm{Ens}_2\times\mathrm{Ens}_3
\end{gathered}\right.
\arrow[u,bend left=75]
& \deleted{\begin{gathered}
\overset{\approx}{\longrightarrow}\unsure{\mathrm{Ens}_2}\times\mathrm{Car}\\
\text{autodualité}\\
\text{de }\mathrm{Ens}_2\times\mathrm{Car}\\
\text{explicite}
\end{gathered}}
\arrow[loop above,leftrightarrow,looseness=3]
\arrow[dr]
&{}\\
\text{remonter à Chap 3} & \sim(\unsure{\text{Hexagones}})&
\text{TSVP}
\end{tikzcd}}

\margin{\diagram{\begin{tikzcd}[row sep=small,column sep=small]
&\approxeq\\
\text{compos}\arrow[ur]&\mathrm{Car}\arrow[u,"\deleted{\unsure{2}}"']
\end{tikzcd}}}

\diagram{\begin{tikzcd}
\alpha\times f &(\alpha,f)\arrow[l]
\end{tikzcd}}

\deleted{\unsure{5)}}
